\subsection{Uniform finite presentation}

Since virtually nilpotent groups are finitely presented, it is a consequence of Gromov's theorem that every group of polynomial volume growth is finitely presented. The purpose of this paper is to prove a \emph{uniform} version of this fact, stating roughly that every group of polynomial growth has a bounded number of scales witnessing a new relation after the first scale that polynomial growth is witnessed. This result is used in our work \cite{Locality} as part of our proof  of Schramm's locality conjecture for Bernoulli bond percolation \cite{MR2773031}, where it plays an important part in our ``uniformization'' of the methods of Contreras, Martineau, and Tassion \cite{contreras2021supercritical}.  A comparison of our results with  the previous literature is given at the end of this section. 

\medskip

Let us now state our result formally.
Let $G$ be a group with finite generating set $S$, so that $G\cong F_S / R$ for some normal subgroup $R$ of $F_S$. For each $n \geq 0$, let $R_n$ be the set of words of length at most $2^n$ in the free group $F_S$ that are equal to the identity in $G$, and let $\llangle R_n \rrangle$ be the normal subgroup of $F_S$ generated by $R_n$, so that the quotient map $F_S/\llangle R_n\rrangle \to G$ induces a covering map of the associated Cayley graphs that has injectivity radius at least $2^{n-1}-1$ (see \cref{lem:small_presentation_means_balls_look_similar}). We say that $(G,S)$ has a \textbf{new relation on scale n} if $\llangle R_{n+1} \rrangle \neq \llangle R_{n} \rrangle$. A finitely generated group $G$ is finitely presented if and only if it has a new relation on at most finitely many scales, so that the following theorem can indeed be thought of as stating that groups of polynomial growth are ``uniformly finitely presented".

\begin{thm}
\label{thm:main}
For each $K,k<\infty$ there exist constants $r_0=r_0(K)$ and $C=C(K,k)$ such that if $G$ is a group and $S$ is a finite generating set for $G$ with $|S|\leq k$ whose growth function $\operatorname{Gr}$ satisfies $\operatorname{Gr}(3r)\leq K \operatorname{Gr}(r)$ for some integer $r\geq r_0$ then
\[
\#\Bigl\{n\in \mathbb{N} : n\geq \log_2(r) \text{ and $(G,S)$ has a new relation on scale $n$} \Bigr\}\leq C.
\]
 %
\end{thm}

\begin{remark}
Considering the abelian group $\prod_{i=1}^k(\Z/n_i \Z)$ with its standard generating set, where $n_1,\ldots,n_k$ are arbitrary, we see that it is not possible to control on \emph{which} scales we find a new relation; we only claim that the total \emph{number} of scales on which we find a new relation is bounded.
\end{remark}

We will prove the following theorem about the number of times we find an ``unexpected element'' during a breadth-first exploration of a (not necessarily normal) subgroup of a group of polynomial growth; we will see in \cref{sec:wrap_up} that this theorem easily implies \cref{thm:main}.

\begin{thm}[Breadth-first exploration of subgroups]
\label{thm:main_generalized}
For each $K$ and $k$ there exist constants $r_0=r_0(K)$ and $C=C(K,k)$ such that the following holds.
Let $G$ be a group with finite generating set $S$ satisfying $|S|\leq k$, let $H$ be a subgroup of $G$, and for each $n\geq 1$ let $H_n$ be the subgroup of $H$ generated by elements that have word length at most $2^n$ in $(G,S)$. If $r\geq r_0$ is such that $\Gr(3r)\leq K \Gr(r)$ then 
\[
\#\{n \geq \log_2 r :H_{n+1}\neq H_n\} \leq C.
\]
\end{thm}

\medskip

\begin{remark}
We believe that it should be possible to take the constants in \cref{thm:main,thm:main_generalized} to be independent of the size of the generating set. We do not pursue this here.
\end{remark}
